\documentclass[a4paper,12pt]{ltjsarticle}
\usepackage[top=25mm,bottom=25mm,left=25mm,right=25mm]{geometry}
\usepackage{graphicx}
\usepackage{float}
\usepackage{booktabs}
\usepackage{array}
\usepackage{enumitem}
\usepackage{hyperref}
\usepackage{caption}

\begin{document}

\begin{center}
  {\Large \textbf{第4回議事録}}\\[8pt]
  {\large チーム4}
\end{center}

\vspace{4mm}

%-------------------------------------------------------------------
\section*{1.基本情報}
\begin{tabular}{ll}
  \toprule
  日時     & 2026年5月13日（水） \\
  会議場所 & 731教室 \\
  目的     & プロジェクトの計画書について \\
  チーム番号 & 4 \\
  議事録作成者 & 酒井睦基 \\
  \bottomrule
\end{tabular}

\vspace{4mm}
\noindent\textbf{参加者}

\begin{tabular}{ll}
  2年 & 山口汐里，佐藤夏美，成毛海人 \\
  3年 & 細澤悠真，武本龍，酒井睦基 \\
\end{tabular}

%-------------------------------------------------------------------
\section*{2.議題一覧}
\begin{enumerate}[label=議題\arabic*：]
  \item プロジェクトの計画書の原案
  \item スライド作成の役割分担
  \item 個人毎の次回までに実施する内容
\end{enumerate}

%-------------------------------------------------------------------
\section*{3.議題ごとの内容}

\subsection*{議題①：プロジェクトの計画書の原案}
\subsubsection*{決定事項}
・スレーブArduinoはそれぞれ別のPCに接続（シリアル通信の混線を防ぐため）
・システムの開始トリガをロードセルに決定
・ベビーメリーの回転に連続回転サーボモータを使用する（1回転の検知には時間を用いる）
\subsubsection*{議論内容}
音楽を開始するためのトリガと音楽が始まるまでのエンタメ性，音楽開始の流れをどうするのか？ \\
→ロードセルは荷重により抵抗が変更され，HX711モジュールにより増幅された値がArduinoに送信される． \\
この数値が閾値を判断基準とし，ベビーメリーが回転する．その後1回転したら音楽が開始される．
1回転の検知には回転速度から逆算した時間を用いる．

%-------------------------------------------------------------------
\subsection*{議題②：スライド制作の役割分担}

\subsubsection*{決定事項}
技術的・社会的課題：武本\\
提案システムのシステム要件，全体構成，構成要素の説明：2年生\\
スライドの内訳：\\
  システム要件：山口\\
  全体構成：佐藤\\
  構成要素：成毛\\
到達目標：酒井睦基\\
役割分担・機材・スケジュール：細澤\\
まとめ：3年生の誰かが作成
（すべてのスライドが揃ってからでない場合作成不可であるため）
\subsubsection*{議論内容}
社会的背景は武本の調べた内容が採用されたためスライドも技術的・社会的課題の説明を武本に任せる．
システム全体の流れの再認識と理解を深めてもらうために提案システムのスライド作成を2年生に任せる．
提案システムのスライド作成の分担，システム要件，全体構成，構成要素については当人たちに話し合ってもらった．
その結果システム要件を山口，全体構成を佐藤，構成要素を成毛が担当することとなった．
到達目標に関しては論文調査を行っている酒井に任せる．
残った役割分担・機材・スケジュールのスライド作成を細澤に任せる．
なお，まとめのスライドに関しては他のスライドがすべて揃ったタイミングで統合する3年生が担当する．

\subsection*{議題③：個人毎の次回までに実施する内容}

\subsubsection*{決定事項}
2年生
・土曜日までにスライドver.1を3年生に提出
・発表練習の動画を送信しそれに対する質疑応答リストを作成

3年生
・提出してもらったスライドver.1を修正
・月曜日に発表練習の動画を撮り2年生に提出
・質疑応答リストを確認し回答とスライドの最終調整を行う

\subsubsection*{議論内容}
スライドの修正期間は複数設ける必要がある．\\
→土曜日を最初の期限とし，そこで3年生がスライドの確認と修正，修正したスライドに関しては2年生に送信する．
スライドの修正と統合は3年生が行う．\\
スライドの発表が必要．\\
→月曜日に3年生で集まって発表練習を行う．
また練習の動画を撮り2年生に送る．2年生には月曜日中にその動画を確認してもらい，質疑応答をリスト化してもらう．
3年生はそのリストを見てスライドの修正と質問に対する回答を考えておく．

\section*{4.次回予定}
\begin{tabular}{ll}
  \toprule
  日時・場所 & 2026年5月20日（水）・731教室 \\
  議題       &  計画発表 \\
  その他，特記事項 & なし \\
  \bottomrule
\end{tabular}

\end{document}
